\section{Related Work}
\label{sec:s2-related}

The communication decision in networked coordination contains three distinct
questions: whether the state is informative enough to transmit, what the
sender causally knows about receiver freshness, and how a feedback action
changes contention on the shared medium. Existing work usually resolves one
or two of these questions. We therefore organize the literature by decision
information and network mechanism rather than list papers by application.

\subsection{State relevance: event-triggered coordination}

Classical event-triggered control replaces periodic updates with transmissions
driven by a local state or estimation error
\cite{astrom2002riemann,tabuada2007event,heemels2012introduction}. Distributed
extensions establish consensus and coordination guarantees, avoid Zeno
behavior, and introduce dynamic thresholds
\cite{dimarogonas2012distributed,seyboth2013eventbased,girard2015dynamic,yi2017distributed,nowzari2019eventtriggered}.
This line extends to multi-UAV formation under directed topologies, delays,
disturbances, and dynamic triggers
\cite{yin2023eventbased,ji2023dynamic,chen2024distributed,yang2025fencing},
including finite-time event-triggered experiments with Bebop2 quadrotors
\cite{wang2022finiteuav}. Event-triggered formation is therefore not the gap.

Loss-aware event triggering is also established. Wang and Lemmon study
distributed event triggering with delay and packet loss
\cite{wang2011eventtriggering}; Dolk and Heemels explicitly distinguish ETC
with and without acknowledgements \cite{dolk2017packetlosses}; Garcia
\emph{et al.} address non-consistent broadcast loss and delay while documenting
earlier receiver-specific ACK/retransmission protocols
\cite{garcia2016unreliable}; and Viel \emph{et al.} trigger formation updates
using estimates of an agent's own state as evaluated by its neighbors under
packet loss \cite{viel2022packetlossformation}. ACK-aware ETC and
receiver-dependent versions of a broadcast state are consequently foundations,
not contributions of this paper.

One remaining abstraction is receiver observability. A sender often resets its
broadcast error when it transmits, effectively treating the transmitted sample
as the receivers' new reference despite queueing or forward loss. Kesper
\emph{et al.} attach age timers to last-broadcast states in learned distributed
event-triggered control \cite{kesper2023toward}, but the timer is not a delayed
receiver-specific acknowledgement of an accepted generation. State relevance
answers whether the plant changed, but not whether a particular neighbor
received the change.

\subsection{Time and semantic relevance}

AoI measures elapsed time since generation of the newest update at a receiver
\cite{kaul2012realtime,sun2017update,yates2021aoisurvey}. AoI-aware networked
control schedules scarce transmissions or constrains estimation quality.
Mamduhi \emph{et al.} combine age- and event-based scheduling across multiple
feedback loops \cite{mamduhi2020freshness}; Wang \emph{et al.} impose
freshness constraints in event-triggered Kalman consensus
\cite{wang2021freshness}; and Lin \emph{et al.} integrate AoI into
event-triggered load-frequency control \cite{lin2023eventtriggered}. These
works establish that freshness is control relevant and that pure state
innovation can leave a receiver stale.

Age is not identical to task value. AoI can underperform state- or
value-aware policies in networked control \cite{ayan2019aoivoi,mamduhi2020freshness}.
Semantic metrics make the same distinction: Age of Incorrect Information
penalizes information that is both stale and wrong \cite{maatouk2020aoii},
whereas Age of Changed Information combines elapsed time with content change
\cite{wang2022aoci}. Decentralized AoII has also been examined under
slotted-ALOHA collisions \cite{nayak2023decentralizedaoii}. Hence neither
content-aware freshness nor collision-aware age scheduling alone is the
present contribution. The unresolved issue is how such relevance is inferred
causally for every formation neighbor and translated into actual shared-medium
traffic.

\subsection{ACKs and delayed two-way feedback}

With unreliable delivery, sender transmission history does not reveal
receiver AoI. Ceran \emph{et al.} optimize ACK-driven sampling and
retransmission under resource constraints \cite{ceran2019average}. Pan
\emph{et al.} derive threshold sampling with random two-way delay and ACK
feedback \cite{pan2023twoway}. Tahir \emph{et al.} consider decentralized
sensors sharing capacity-limited non-FIFO duplex channels; each sensor uses a
particle filter to maintain a belief over receiver AoI under delayed ACKs
\cite{tahir2024collaborative}. Onozuka \emph{et al.} event-trigger forward and
feedback traffic in railway control \cite{onozuka2024aoi}.

This strand resolves the causal observability problem omitted by classical
event triggers, and the delayed-ACK belief method is consequently implemented
as an experimental baseline rather than cited only for positioning. Its main
objective is receiver age or remote estimation. It does not directly answer
whether a continuous formation state is newly informative to a particular
neighbor, or whether a standalone ACK changes future broadcast demand enough
to repay its own airtime and collisions.

That resource question also has direct prior art. Munari and Badia
analytically identify reliability/feedback-cost regions in which requesting
feedback is or is not beneficial for finite-horizon AoI optimization
\cite{munari2025feedback}. Their single-sensor opportunity-budget model does
not include formation semantics, broadcast neighbor beliefs, piggyback versus
standalone contention, or closed-loop DATA suppression. It does mean that the
generic claim ``feedback is sometimes worth its cost'' is not novel. Our gap
is the mediated path from an ACK, through sender belief and future DATA
actions, to receiver freshness and control.

\subsection{Random access and MAC-dependent freshness}

The access mechanism is already known to change AoI. Kadota and Modiano analyze
stochastic-arrival AoI under slotted ALOHA and CSMA
\cite{kadota2021randomaccess}; Mena and N\'u\~nez derive AoI expressions and a
MAC-selection method for CSMA/CA and TDMA in IoT networked control
\cite{mena2023macperspective}; and Peng \emph{et al.} use piggybacked collision
information in an AoI-optimized V2X sidelink MAC
\cite{peng2021piggyback}. Thus neither MAC-dependent AoI nor piggyback signaling
is the contribution. The open issue is whether the \emph{feedback route}
changes the transmitter's semantic update process enough to reverse its
receiver/control value on a shared broadcast medium.

\subsection{Shared-medium UAV networking}

UAV-network research emphasizes mobility, interference, bandwidth limits, and
topology changes \cite{gupta2016survey,zeng2016wireless,mozaffari2019tutorial}.
WiSwarm is particularly relevant: it implements an AoI-based application layer
over Wi-Fi for collaborative UAV teams, including stale-packet handling,
traffic priorities, collision prevention, and flight experiments
\cite{tripathi2023wiswarm}. It demonstrates why logical directed-link counts
are insufficient: broadcast aggregation, queueing and medium access change
delivered freshness and resource cost.

The complementary limitation is that network-oriented systems do not isolate
the present receiver-specific innovation/confirmation mechanism, while many
formation-control papers abstract the MAC. We bridge these layers with an
explicit but radio-agnostic shared-medium kernel. This is stronger than a
broadcast accounting proxy but weaker than a named IEEE PHY/MAC or hardware
experiment; trace and radio validation remain a separate stage.

\subsection{Gap and positioning}

The literature has solved neighboring pieces: semantic state triggers for
coordination, receiver-age scheduling, delayed-ACK receiver beliefs, and
shared-medium UAV networking. The joint gap is not the union of familiar
ingredients. It is the decision mechanism created when receiver-confirmed
semantic triggering, broadcast aggregation, cumulative feedback, finite
contention, and formation control interact.

Our central question is correspondingly conditional: when does faster causal
confirmation improve the receiver-side control process after paying direct and
collision-mediated cost? The answer cannot be inferred from ACK delay alone.
This paper combines pathwise accounting, focal-decision counterfactual replay,
disjoint directional validation, and a preregistered within-seed factorial
interaction to show that the relative route effect reverses with the configured
access mechanism in the declared cell. The contribution is an auditable
causal/accounting framework and a bounded MAC-conditioned route effect, not a priority claim
for event triggering, AoI, ACK inference, broadcast, or CSMA/ALOHA separately.
